\documentclass{beamer}

% You can also use a 16:9 aspect ratio:
%\documentclass[aspectratio=169]{beamer}
\usetheme{TACC16}
\usepackage{graphicx}

% It's possible to move the footer to the right:
%\usetheme[rightfooter]{TACC16}


% Detailed knowledge of application workload characteristics can
% optimize performance of current and future systems. This may sound
% daunting, with many HPC data centers hosting over 2,000 users running
% thousands of applications and millions of jobs per month.  XALT is an
% open source tool developed at the Texas Advanced Computing Center
% (TACC) that collects system usage information to quantitatively report
% how users are using your system. This session will explore the
% benefits of detailed application workload profiling and how the XALT
% tool has helped leading supercomputing sites unlock the power of their
% application usage data. 


\begin{document}
\title[XALT]{XALT Monthly Meeting: Controlling Executables Tracked}
\author{Robert McLay} 
\date{Sept.  9, 2021} 

% page 1
\frame{\titlepage} 

% page 2
\begin{frame}{XALT: Outline}
  \center{\includegraphics[width=.9\textwidth]{XALT_Sticker.png}}
  \begin{itemize}
    \item XALT's DB can easily get overwhelmed
    \item Three ways to control what gets recorded
    \item TACC's goals are to track science
    \item Config/TACC\_config.py discussion
    \item Sampling
  \end{itemize}
\end{frame}

% page 3
\begin{frame}{Preventing XALT's DB getting overwhelmed}
  \begin{itemize}
    \item Must do something to prevent too much data
    \item Want to understand what is run on your system w/o tracking everything
    \item Tales of Woe:
      \begin{itemize}
        \item One 2 hour job ran over 2M times $\Rightarrow$ 4.5 days
          of data
        \item One user used a 4 core MPI program to train a neural network
      \end{itemize}
    \item XALT supports sampling for both MPI and Non-MPI programs 
  \end{itemize}
\end{frame}

% page 4
\begin{frame}{Three ways to control the amount of recorded data}
  \begin{itemize}
    \item Use path\_patterns to ignore many executables based on path.
    \item Sampling to reduce the amount of executables tracked based
      on runtime.
    \item Pre-ingest filtering to sample data already tracked.
      (only works w/ storing data in the MySQL DB)
  \end{itemize}
\end{frame}

% page 5
\begin{frame}{Tracking Science}
  \begin{itemize}
    \item We want to know what science is being done
    \item What apps are run and build by users or in /opt/apps
    \item Sample apps like cp, perl, awk, and sed
    \item Sites are free to configure XALT as they see fit.
  \end{itemize}
\end{frame}

% page 6
\begin{frame}{Config.py $\Rightarrow$ generated code}
  \begin{itemize}
    \item For speed reasons, the patterns are converted to a lex program
    \item This lex program is converted into executable code
    \item These patterns are then directly callable from the XALT
      shared library.
    \item This pattern matching is very fast.
    \item Downside: Any changes in path\_patterns must be rebuilt.
  \end{itemize}
\end{frame}

% page 7
\begin{frame}{Config/TACC\_config.py}
  \begin{itemize}
    \item This file is what TACC uses.
    \item Note that the first match found is used.
    \item Mark R, MATLAB and Python as PKGS that track imports
    \item Keep cp, sed, awk, perl and similar from /bin or /usr/bin
    \item Ignore anything else in /usr/bin or /bin
    \item Ignore compilers, mpi wrappers, shell scripts (/bin/bash is
      ignored)
    \item Ignore test programs from Cmake and autoconf's config
  \end{itemize}
\end{frame}

% page 8
\begin{frame}[fragile]
    \frametitle{Sampling}
 {\small
    \begin{semiverbatim}
interval\_array = [
    [    0.0,              0.0001 ],
    [ 1800.0,              0.01   ],
    [ 7200.0,              1.0    ],
    [ sys.float\_info.max,  1.0    ]
]
mpi\_interval\_array = [
    [    0.0,              0.0001 ],
    [  900.0,              0.01   ],
    [ 1800.0,              1.0    ],
    [ sys.float\_info.max,  1.0    ]
]
    \end{semiverbatim}
}
\end{frame}

% page 9
\begin{frame}[fragile]
    \frametitle{path\_patterns}
 {\tiny
    \begin{semiverbatim}
path_patterns = [
    ['PKGS',  r'.*\\/python[0-9.]*'],
    ['KEEP',  r'^\\/bin\\/gawk'],
    ['KEEP',  r'^\\/bin\\/sed'],
    ['KEEP',  r'^\\/bin\\/perl'],
    ['SKIP',  r'^\\/bin\\/.*'],
    ['SKIP',  r'^\\/opt\\/intel\\/.*'],
    ['SKIP',  r'^\\/opt\\/apps\\/git\\/.*'],
    ['SKIP',  r'^\\/opt\\/apps\\/cmake\\/.*'],
    ['SKIP',  r'^\\/opt\\/apps\\/autotools\\/.*'],
    ['SKIP',  r'^\\/opt\\/apps\\/intel[0-9][0-9\\_]*\\/impi\\/.*'],
    ['SKIP',  r'^\\/opt\\/apps\\/gcc[0-9][0-9\\_]*\\/mvapich2\\/.*'],
    ['SKIP',  r'^\\/home1\\/apps\\/intel\\/.*'],
    ['SKIP',  r'.*\\/mpicc'],
    ['SKIP',  r'.*\\/mpicxx'],
    ['SKIP',  r'.*\\/mpif90'],
    ['SKIP',  r'.*\\/conftest'],
    ['SKIP',  r'.*\\/CMakeTmp\\/cmTryCompileExec[0-9][0-9]*'],
    ['SKIP',  r'.*\\/CMakeTmp\\/cmTC_[a-f0-9][a-f0-9]*'],
  ]
    \end{semiverbatim}
}
\end{frame}

% page 10
\begin{frame}[fragile]
    \frametitle{Before XALT-2.10.30}
A site did something like this:
 {\small
    \begin{semiverbatim}
path_patterns = [
    ['KEEP',  r'^\\/opt\\/apps\\/.*'],
    ['SKIP',  r'.*'],
  ]
    \end{semiverbatim}
}
Where XALT was installed /opt/apps/xalt
\end{frame}

% page 11
\begin{frame}{The Config file $\Rightarrow$ compiled code}
  \begin{itemize}
    \item XALT uses python scripts to generate *.lex code for regex processing
    \item It also takes the same files to generate *.h file for
      reporting configuration.
    \item xalt\_configuration\_report is there for users/admins to know
      how XALT is installed.
  \end{itemize}
\end{frame}

% page 12
\begin{frame}[fragile]
    \frametitle{xalt\_configuration\_report (XALT-2.10.30) and later} 
 {\tiny
    \begin{semiverbatim}
*----------------------*
 Array: pathPatternA
*----------------------*
================== src/tmpl/xalt_config.py ==================
   0: SKIP => .*\\/bin\\/my\\_uuidgen
   1: SKIP => .*\\/logger
   2: SKIP => .*\\/xalt\\_print\\_os
   3: SKIP => .*\\/xalt\\_syslog\\_to\\_db
   4: SKIP => .*\\/xalt\\_extract\\_record.x
   5: SKIP => .*\\/xalt\\_configuration\\_report.x
   6: SKIP => .*\\/xalt\\_syshost
   7: SKIP => .*\\/xalt\\_record\\_pkg
   8: SKIP => ^\\/opt\\/apps\\/xalt\\/xalt\\/.*
================== /opt/apps/xalt/Config/simpleConfig.py ==================
   9: KEEP => ^\\/opt\\/apps\\/.*
  10  SKIP => .*
================== src/tmpl/xalt\\_config.py ==================
  11: KEEP => .*
    \end{semiverbatim}
}
\end{frame}

% page 13
\begin{frame}{Conclusions}
  \begin{itemize}
    \item Thanks for listening
    \item Questions?
  \end{itemize}
\end{frame}

%
\begin{frame}{License}
\footnotesize
\textcopyright The University of Texas at Austin, \the\year\\
\vspace{0.5 cm}
This work is licensed under the Creative Commons Attribution Non-Commercial 3.0 Unported License. To view a copy of this license, 
visit http://creativecommons.org/licenses/by-nc/3.0/\\
\vspace{0.5 cm}
When attributing this work, please use the following text: 
"Designing and Administering Large-scale Systems", Texas Advanced Computing Center, \the\year. Available under a Creative Commons Attribution Non-Commercial 3.0 Unported License.
\newline
\newline
\includegraphics[width=1cm]{./figures/cc.png}
\end{frame}

\end{document}
